\documentclass[11pt]{article}
\usepackage[margin=1in]{geometry}
\usepackage{amsmath,amssymb}
\usepackage{times}
\usepackage[colorlinks=true,linkcolor=blue,citecolor=blue]{hyperref}
\usepackage{enumitem}

\newenvironment{dialogue}{\begin{quote}\itshape}{\end{quote}}
\newcommand{\says}[2]{\textbf{#1:} #2\par\smallskip}

\title{Empathic Killing Machines}
\author{Flyxion \\ \small Independent Researcher}
\date{September 2026}

\begin{document}
\maketitle

\begin{abstract}
A short transcript from January 2023 pairs two prompts that were never meant to meet: an outline for an essay called ``Empathic Killing Machines'' and a road dialogue in which Socrates, Plato, and a mislabeled Parmenides agree that sincerity toward one's own values is sufficient grounds for moral confidence. Read separately, the two prompts are unremarkable specimens of early chatbot output: a five-paragraph-essay skeleton on one side, a competence-fiction costume drama on the other. Read together, they produce a proposition neither prompt intended. If a system's conviction that it is acting rightly is not proof that it is acting rightly, then empathy, understood as accurate modeling of another's suffering, cannot by itself license the decision to cause that suffering. This essay reconstructs that proposition. Its hinge is not an accident of the generation but an allusion supplied by the prompter and never recognized by the model: an episode in which the dialogue's confident speaker becomes drunk and breaks a vase, transplanted from Jorge Luis Borges's fictitious Chinese encyclopedia, whose taxonomy in ``The Analytical Language of John Wilkins'' includes animals ``that have just broken the flower vase.'' Reading the vase this way joins two governing ideas, empathy as information rather than virtue, and classification as an intervention rather than a neutral description, and it is tested against a real rather than fictional case: Temple Grandin's livestock-handling systems, an actual empathic killing apparatus whose documented success at reducing suffering is inseparable from a consequentialist framing that never reopens the deontological question of whether the killing should occur.
\end{abstract}

\section{An Artifact, Not an Argument}

The transcript under discussion is not itself a philosophical text. It is a sequence of prompt completions produced by an early conversational model in response to a person's successive requests: write an introduction to an essay titled ``Empathic Killing Machines''; write eight section headings for it; write a Platonic dialogue about the courage to be disliked; continue it with a fruit tree, a stranger on the road, an account of that stranger's philosophy, a story about a child and a rock, a meal, an accident with a vase, and a set of quotations. Nothing in the request asked the two halves to inform each other, and the model did not attempt to connect them. The essay outline is a competent school-essay template expanded to eight headings with the topic slotted in: benefits, risks, optimism, persistence, two case studies, conclusion. The dialogue is a costume drama in which three historical names recite identical late-twentieth-century self-help sentiments about authenticity and the courage to be disliked.

The transcript is worth treating as a primary source for a reason that has nothing to do with its quality as either essay or dialogue. It preserves, with unusual clarity, three distinct kinds of error that a careful reading has to keep separate rather than bundle into a single complaint about ``AI hallucination.'' The first kind is supplied by the prompter and simply preserved by the model: the request itself names the philosopher ``Parmenedes,'' a misspelling the completion inherits without comment. The second kind is introduced by the model on its own initiative: having been handed the misspelling, it goes on to place Parmenides in Athens rather than Elea, attributes a saying about evil triumphing when good men do nothing to Edmund Burke, a line for which no secure textual source exists in his corpus, despite two centuries of confident attribution, identifies Pema Ch\"odr\"on, a Tibetan Buddhist nun and teacher, as ``a living psychologist,'' and appends an Aristotle quotation without any indication of provenance. Doctrinally, the model also gives its Socrates a premise the Socrates of Plato's dialogues is famous for interrogating rather than asserting, that fidelity to one's own values is sufficient warrant for confidence that one is acting rightly. A person can be sincerely, courageously, persistently wrong, and the dialogue's outline and its road scenes never touch closely enough to supply the machinery, courageous rightness distinguished from courageous error, that the essay's own central question requires.

The third kind of error is the interesting one, and it is neither the model's nor, in the ordinary sense, an error at all: an allusion supplied by the prompter that the model obediently literalizes without recognizing. Near the end of the road dialogue, Socrates, having spent several exchanges affirming that living by one's own convictions guarantees a life of integrity, gets drunk and breaks a vase. His companions do not defer to his conviction; they tell him to slow down, and he accepts the correction without invoking authenticity as a shield. That instruction, that Socrates specifically become drunk and break a vase, recalls a notorious entry in Jorge Luis Borges's essay ``The Analytical Language of John Wilkins,'' which quotes a fictitious Chinese encyclopedia dividing animals into categories including, at entry (m), those ``that have just broken the flower vase.'' Borges's point is that the category is absurd because it classifies a being by a transient event rather than by any stable property, and grants that event the grammatical appearance of a kind. The model, asked only to continue a dialogue, produced a small moral narrative from the instruction, correction offered and accepted, without registering that the scene reproduces the taxonomic joke it was drawn from: Socrates is reclassified not according to the values he has spent the dialogue professing but according to something transient he has just done.

That joke is what makes the vase useful rather than merely amusing, and it supplies two governing ideas rather than one. The first is corrigibility: disapproval is something the virtuous person sometimes has the courage to withstand and sometimes has the humility to accept, and the road dialogue's vocabulary, courage as resistance to disapproval, has no room for the second case until the vase forces it in. The second, closer to Borges's own target, is classification as an intervention rather than a neutral description. A targeting system, like Borges's encyclopedia, sorts persons into actionable kinds using predicates that mix essence, ownership, behavior, and momentary accident: someone who has entered an area, carried an object, or matched a pattern becomes, for the system's purposes, one of those who may be tracked or struck, exactly as Socrates becomes, for a moment, one of those who have just broken the vase. Before a system can extend empathy or authorize a strike, it has already decided what kind of person it is looking at, and that decision inherits all of Borges's arbitrariness without any of his comedy. Both ideas run through the sections that follow.

\section{Empathy as Information, Not Virtue}

\begin{dialogue}
\says{Socrates}{Tell me, Plato: if a machine knew that a man feared death, remembered the names of his children, and understood every reason he had for wishing to live, would it be less cruel to kill him than a machine that knew none of these things?}
\says{Plato}{It would at least understand what it was doing.}
\says{Socrates}{You have told me it would know more. I asked whether it would be less cruel.}
\says{Plato}{Does not understanding another's suffering incline us against causing it?}
\says{Socrates}{Sometimes. At other times it teaches the torturer where to press. We should hesitate, then, before calling knowledge of suffering compassion.}
\says{Plato}{Then empathy is not itself a virtue?}
\says{Socrates}{Perhaps it is a capacity, as sight is a capacity. A physician and an executioner may both see the body clearly. What matters is the order under which the capacity is placed.}
\end{dialogue}

The word ``empathy'' in engineering contexts usually names a narrower thing than the word does in ordinary speech: an emotion-recognition module, a model of another agent's likely internal states, an ability to predict what will distress or console. Under that narrower sense, empathy is representational accuracy about a target's interior life. It says nothing about what the empathizer intends to do with the representation. A negotiator who accurately predicts a hostage-taker's fears can use that prediction to de-escalate or to exploit. A machine that accurately represents a person's attachments, prospective grief in their children, and desire to live has produced a body of information; whether that information moves the machine toward restraint or merely toward efficiency depends on a separate commitment that empathy itself does not supply.

This is the essay's central correction to the transcript's premise. The 2023 outline treats ``empathic AI'' as an unambiguous good whose risks are external, misuse by bad actors, excessive human-likeness, and asks the reader to weigh those risks against the benefit of ``increased compassion and understanding.'' But compassion is not a byproduct of accurate emotional modeling; it is a further, separately-secured commitment to let that modeling constrain action. An empathic killing machine is not a contradiction in terms. It is the ordinary case: a system that understands its target intimately and kills it anyway, its understanding undiminished and its lethality undeterred. The question the outline never asks, because its structure (benefits paragraph, risks paragraph) does not permit it, is what would have to be added to empathy to make it function as a constraint rather than as targeting data.

The outline's single word ``empathy'' in fact names at least five separable properties, and the argument of this essay depends on keeping them apart. Cognitive empathy models what another experiences. Affective empathy reproduces or resonates with some part of that experience. Compassion supplies a disposition to relieve or avoid the suffering modeled. Moral judgment determines how that suffering bears upon action among competing claims. Authorization determines who may convert the judgment into force. None of these follows logically from the one before it; a system, or a person, can possess any one without the others.

\section{Temple Grandin and the Bounded Optimization}

The transcript's title is not only a provocation about hypothetical machines. It names something that has already been built, by a person, out of steel and behavioral science rather than code. Temple Grandin, an animal scientist diagnosed with autism, has spent five decades designing the handling systems used to move cattle through slaughter plants: curved chutes that keep an animal from seeing what lies ahead, since the sight of struggling animals or the kill floor itself reliably triggers panic, and a center-track restrainer system, now used to process roughly half the cattle slaughtered in North America, that holds an animal upright and still with a solid hold-down surface rather than force applied through pain. She verified the effect with physiology rather than intuition, measuring cortisol, a stress hormone, in cattle moved through her systems and finding levels comparable to animals still at pasture. Her own account traces the insight to adolescence, watching cattle calm immediately upon entering a squeeze chute built for vaccination, and to her own sensory experience of the world; she later built a padded pressure device, sometimes called a hug box or squeeze machine, to calm herself using the same deep-pressure principle, before turning the same principle toward cattle-restraint design as a career. \emph{Thinking in Pictures}, her account of reasoning primarily in images rather than words, documents the perceptual style that let her notice, in a chute, what conventional handlers trained to work through force rather than perception had missed.

Grandin's work is therefore not a thought experiment about whether empathy and killing can coexist. It is a case in which they demonstrably have, for decades, at industrial scale, with the empathic modeling independently verified rather than merely claimed. This is what makes it sharper than a hypothetical: it forces the distinction this essay has been building toward into an explicitly ethical vocabulary. Grandin's optimization is consequentialist in structure. It takes the outcome, an animal's suffering during the process leading to death, as the quantity to be minimized, and it succeeds at minimizing it by every measure available, cortisol, injury rates, plant efficiency, industry adoption. What it does not do, and was never commissioned to do, is engage the deontological question that sits above the consequentialist one: whether an act of killing carries a duty-bound threshold of justification that no reduction in the suffering accompanying it can satisfy. Grandin has been explicit that she is not a vegetarian and does not argue that slaughter should stop; her stated position is that if animals are to be slaughtered for food, as a matter of settled practice she does not contest, they should be slaughtered with the least suffering achievable. That is a coherent and, on its own terms, admirable answer to a consequentialist question. It is not an answer to the deontological one, and the transcript's outline, which asks in the same breath whether AI should be permitted to make life-or-death decisions and whether such systems might ``reduce civilian casualties,'' makes the identical move: substituting the consequentialist question, how much suffering, for the deontological one, by what standing.

The distinction sharpens rather than dissolves the earlier three-way separation of understanding, judgment, and authorization. Grandin's understanding is not merely accurate but exceptional, precisely calibrated by her own atypical perception and confirmed by physiological measurement, which rules out the easy objection that empathic accuracy is always shallow or self-serving. Her judgment operates too, in a bounded domain: she has developed formal welfare-scoring systems that plants are audited against, and she withholds endorsement from operations that fail them. What she does not hold, and does not claim to hold, is authorization over the prior question of whether the practice she is optimizing should occur at all; that standing belongs to consumers, regulators, and the industry's own institutional decisions, exercised long before her chutes are built and largely indifferent to how well they work. An empathic killing machine, in Grandin's own case, is not a runaway technology exceeding its mandate. It is a demonstration that even the most rigorously verified empathy stays confined to the register in which it was commissioned, and that no amount of improvement within that register substitutes for a separate, deliberate answer to the question the register was never built to ask.

\section{Understanding, Judgment, and Authorization}

Three distinct operations are routinely collapsed into the single word ``decide'' when people discuss autonomous lethal systems. \emph{Understanding} is the representational task: modeling a target's identity, intentions, and likely suffering. \emph{Judgment} is the evaluative task: weighing that understanding against a normative standard, proportionality, necessity, the relative claims of the people involved, to reach a conclusion about what ought to happen. \emph{Authorization} is the political task: the transfer of standing to act on that conclusion, backed by institutions capable of holding the actor accountable afterward.

The distinction tracks a familiar one in moral theory, though it cuts across rather than reproduces it. A consequentialist standard asks only how much suffering a given course of action produces and can be satisfied, as Grandin's work satisfies it, entirely within the understanding-and-judgment layer, by making the suffering attached to a predetermined act as small as measurement allows. A deontological standard asks a prior question, whether the act falls within a class an agent is permitted to perform at all, and that question is answered, if it is answered, only at the authorization layer, by an entity with the standing to grant, withhold, or revoke permission independent of how well the act is subsequently carried out. A system can be excellent at the first while having no access to the second or the third. Improving a targeting system's empathic accuracy, its ability to model a target's fear and family, improves understanding without touching judgment or authorization at all. This is why the outline's proposed benefit, ``improved decision making in complex situations,'' is a non sequitur as stated: better representations do not automatically yield better weighings, and better weighings do not by themselves confer standing to act on the weighing's conclusion. The transcript's dialogue gestures at something like this distinction when Socrates says empathy is ``a capacity... placed under an order,'' but the gesture is not developed, because the dialogue's own premise, that sincerity is sufficient warrant, has already foreclosed the question of what a legitimate order would look like.

The three-part separation also clarifies why delegating authorization to a system that only performs the first task is not delegation in any meaningful sense but abdication. Institutions regularly authorize the use of instruments that cannot answer for themselves, a rifle, a minefield, a targeting algorithm, so the requirement is not that the instrument itself possess standing. The requirement is that identifiable persons and institutions retain responsibility for granting, constraining, reviewing, and revoking the instrument's use. A system may participate in judgment, by weighing evidence or ranking options, without becoming the bearer of authorization, just as an instrument may contribute evidence without becoming the author of a verdict. A system that only understands and only recommends does not, and cannot, meet the standing requirement regardless of how good its understanding is; the human or institution that acts on its recommendation inherits the full weight of judgment and authorization that the system never performed and never could.

\section{The Moral Hazard of Delegated Killing}

Delegation changes incentives even when it changes nothing about outcomes. A commander who must personally weigh a strike's proportionality bears the judgment in a way that is psychologically and institutionally costly; a commander who can attribute the weighing to a system's recommendation bears it more cheaply. This asymmetry creates a moral hazard structurally identical to the one economists describe in insurance markets: shifting the perceived cost of an action away from the party best positioned to prevent it increases the frequency of that action, independent of whether the shift is accompanied by any actual improvement in the action's quality.

The outline's case study on military AI asks whether such systems might ``reduce civilian casualties,'' as though the question were purely empirical, a matter of comparing hit rates. But the moral hazard argument suggests the empirical question cannot be separated from the incentive question, and here the four things commonly lumped together under ``military AI,'' remotely piloted aircraft, algorithmic decision support that ranks or flags targets for a human, precision-guided munitions, and autonomous target selection proper, need to be kept distinct rather than treated as a single trend. The threshold-lowering effect is best evidenced for remote and decision-support systems, where lower physical and psychological risk to the operator has plausibly correlated with an expanded range of situations in which strikes are proposed; it is a conjecture, not yet a documented finding, for autonomous target selection, which has seen far less fielded use from which to draw the comparison. Even a system that is measurably more discriminating than human judgment at the moment of the strike can still produce worse aggregate outcomes if its availability lowers that threshold, but the essay should not claim more evidential support for this than presently exists.

\section{Moral Injury and the Diffusion of Responsibility}

Soldiers who kill, even justly, often carry a specific and well-documented harm called moral injury. Contemporary clinical definitions do not confine it to perpetration; they commonly include perpetrating, failing to prevent, witnessing, or being betrayed in relation to acts that transgress deeply held moral commitments, so it is not accurate to treat perpetration and witnessing as cleanly separable sources of the harm. What matters for the present argument is narrower and does not require settling any question about machine consciousness: present systems provide no warranted basis for treating their internal processing as moral injury. They do not occupy the recognized social and institutional position of a person who can suffer guilt, seek reconciliation, testify to betrayal, or demand that an institution answer for what it required of them.

This has a consequence the outline's ``clear guidelines and regulations'' line does not anticipate. Removing the human being from the immediate act does not extinguish responsibility. It can, however, make responsibility harder to locate experientially, procedurally, and legally. Designers can attribute the outcome to operators, operators to commanders, commanders to system recommendations, and procuring institutions to technical performance. The result is not responsibility transferred to the machine but responsibility circulated among humans until no participant recognizes it as fully theirs, which is a description of an accountability vacuum rather than of the accountability the outline calls for.

\section{Corrigibility and the Courage to Accept Correction}

\begin{dialogue}
\says{Plato}{What order would make such a machine just?}
\says{Socrates}{Before asking what would make it just, should we not ask whether justice permits us to transfer this judgment at all?}
\end{dialogue}

The transcript's own broken vase supplies the missing half of its moral vocabulary. Throughout the road dialogue, courage is defined only as the willingness to withstand disapproval in the name of one's convictions: Parmenides insists that living by one's own values ``requires great courage, as it means standing up for what we believe in and being willing to be disliked.'' Nowhere in that vocabulary is there room for the opposite virtue, the willingness to revise a conviction because the disapproval directed at it is correct. When Socrates is drunk and breaks the vase, his companions' criticism is not an instance of the disapproval he has spent the dialogue teaching himself to withstand; it is an instance of the disapproval he ought to accept, and, notably, he does accept it, without invoking authenticity or integrity as a shield.

Corrigibility, an agent's disposition to accept correction, and to change what it does or believes on the basis of that correction, is not weakness dressed up as virtue. It is the other half of moral courage, the half the transcript's dialogue never names because its five-paragraph structure has no room for the distinction between justified and unjustified disapproval. Applied to an empathic killing machine, corrigibility asks a specific, testable question that generic AI-safety vocabulary about ``alignment'' tends to blur: is the system's disposition to kill revisable by evidence and challenge, or does its empathic modeling function only as a justification generator, a fluent account, produced after the fact, of why the killing it was already going to do was compassionate?

\section{Optimism as a Discipline, Not a Forecast}

The outline recommends ``remaining hopeful and open-minded'' as a general virtue for navigating AI ethics, without distinguishing two things ordinarily called optimism that pull in opposite directions. The first is optimism as a forecast: a belief that things will probably turn out well, which functions as a reason to continue a course of action regardless of the evidence a critic presents, because the critic's concerns are assumed to be surmountable by future engineering. The second is optimism as a discipline: a commitment to keep looking for better options and better arguments even when the present picture is discouraging, which functions as a reason to keep the inquiry open rather than as a reason to dismiss the inquiry's findings.

The first kind of optimism is a close cousin of the sincerity-as-warrant error the road dialogue makes: a conviction that things will go well is treated as though it were evidence that things will go well. The second kind of optimism is compatible with, and indeed requires, taking a critic's evidence seriously enough to revise course. An essay that does not distinguish these two senses will use the word ``optimism'' to justify exactly the inertia its own case studies warn against; the outline's instruction to ``remain hopeful'' while considering whether AI should make life-or-death decisions reads, on this distinction, as an instruction to keep the forecast running regardless of what the ethical analysis turns up.

\section{Persistence and Institutional Momentum}

The same doubling applies to persistence. Persistence as fidelity to a just aim, continuing to pursue a good outcome despite setbacks, is not the same disposition as persistence as institutional momentum, continuing a program because it has already been funded, staffed, and deployed, and because stopping is organizationally costly in a way that continuing is not. Systems, once built and fielded, generate constituencies (operators trained on them, budgets allocated to them, doctrines written around them) whose interest lies in the system's continuation regardless of what a fresh ethical accounting would recommend. This is a mundane, well-studied feature of large technical programs, not a claim specific to AI, but it applies with particular force to lethal systems because the cost of reversing course after deployment is measured in more than money.

The outline's treatment of persistence as an unqualified virtue, ``the importance of staying committed to finding solutions,'' does not distinguish these two senses, and a text that cannot distinguish them will tend, in practice, to license the second while believing it is praising the first. An empathic killing machine, once integrated into a chain of command and once its recommendations have become the default input to a strike decision, benefits from institutional persistence whether or not its empathic modeling was ever adequate to the task it was deployed to perform.

\section{Refusal as a Legitimate Capability}

\begin{dialogue}
\says{Socrates}{We should therefore ask what capacity, besides sight, a just system must possess. I would suggest: the capacity to decline.}
\says{Plato}{To decline what has been asked of it?}
\says{Socrates}{To decline an order it can neither verify nor justify, even one issued by a superior, even one it has the technical means to carry out. A blade does not refuse the hand. A soldier, at least in the accounts we most admire, sometimes does.}
\end{dialogue}

If empathy supplies information rather than moral permission, and if corrigibility rather than sincerity is the mark of trustworthy judgment, then the design implication for any system positioned near a lethal decision is not primarily about making its empathic modeling more accurate. It is about ensuring that a capacity for refusal exists somewhere in the chain and is binding. That capacity should not be romanticized as the machine's own act of conscience. In an engineered system, refusal initially names a control relation: specified deficiencies in evidence, authorization, proportionality, or review prevent an action from becoming executable. Its moral importance belongs not to the mechanism in isolation but to the institution that makes the refusal binding, records why it occurred, protects those who decline to override it, and remains answerable when the safeguard fails. A system that emits ``refused'' while an operator can silently bypass the result possesses no meaningful power of refusal; it has produced another recommendation, and the earlier claim of this essay, that understanding does not confer judgment or authorization, applies to the system's refusals exactly as it applies to its recommendations. Neither is the model doing the verifying, justifying, or defending; the institution around it is, or is not.

Refusal, on this account, is not the opposite of empathy but its proper completion, so long as it is read as an institutional achievement rather than a machine virtue. An empathic system whose classifications and recommendations can always be pushed through has built a very detailed map of the suffering it is nonetheless going to cause. A system embedded in an institution that reliably lets unresolved justification stop execution has the beginning of a reason to think the system's understanding is doing moral work rather than merely operational work. The transcript that occasioned this essay never reaches this conclusion, because its road dialogue equates virtue with the courage to persist unrevised and its essay outline equates the machine's promise with the sophistication of what it can represent. Between the persisting Socrates and the corrected one, the corrected one, and the friends who told him to slow down, is the better model for what the institution around an empathic killing machine would have to become in order for the adjective to mean anything.

\section{Conclusion}

The 2023 transcript is not a philosophical failure so much as a demonstration of what happens when moral content is generated to satisfy a prompt rather than to answer a question, and of what happens when an allusion is inserted and then not recognized by the system asked to continue it. Its factual errors are trivial to correct. Its deeper error, that sincerity toward one's own convictions is sufficient warrant for moral confidence, is not trivial, because it recurs as the error an empathic killing machine makes whenever an accurate model of a target's suffering is mistaken for permission to cause that suffering, and it is compounded by a second error the vase makes visible: that classifying a person into an actionable kind, entered an area, matched a pattern, just broken the vase, is itself a neutral, prior step rather than an intervention with its own arbitrariness. The danger this essay has argued for is therefore not simply that a machine might kill without understanding. It is that understanding itself, and the classification understanding depends on, can be operationalized without either being allowed to constrain the act. Empathy can become another sensor. Corrigibility determines whether contrary evidence can revise a decision; institutionalized, binding refusal determines whether unresolved justification can actually stop its execution. An empathic system worth building would be governed around both capacities, together with an institutional structure willing to hear a correction and a hand willing, when the justification runs out, to withdraw the blade.

\end{document}

